\section{Introduction}
\label{sec:intro}

An offline reinforcement-learning (RL) learner must act from a fixed log; it cannot gather new evidence while it trains~\cite{levine2020offline}. A practitioner who holds a small pilot log therefore faces a planning question before committing to a large collection or training grid: \emph{how much logged data will this learner need before it reliably reaches a target return?} Dataset size alone does not answer it. Two logs of equal size can differ in which state--action pairs they cover, in the quality of the logging behavior and in how much transition noise must be averaged out, and the answer depends on the learner: behavior cloning (BC) cannot exceed a poor behavior, whereas a value-based learner can improve on it if the relevant actions are logged~\cite{kumar2022prefer}.

Three obvious shortcuts are insufficient. Dataset-quality scores~\cite{schweighofer2022dataset,swazinna2021quality,asadulaev2026bwd} summarize a log in one number that is not tied to a learner or a budget. Sample-complexity bounds~\cite{li2024sample,wang2022gapoffline} are worst-case or depend on unknown instance quantities. Training the learner on prefixes of the pilot and extrapolating its learning curve~\cite{domhan2015speeding,adriaensen2023efficient} sees only budgets smaller than the pilot, and its success must itself be estimated off-policy~\cite{le2019batch,voloshin2021empirical}. We did not find, in a scoped search, a method whose output is the \emph{distribution of the first data budget at which a specified offline learner attains sustained success}, estimated from a pilot and scored against censored outcomes.

We study this task in finite stochastic Markov decision processes (MDPs), where every policy can be evaluated exactly, so the target is measured rather than estimated. A forecaster sees one pilot of $10|S||A|$ transitions and public metadata; for BC, conservative fitted Q-iteration (CFQI) or tabular implicit Q-learning (IQL)~\cite{kostrikov2022iql} it outputs seven probabilities over the budgets $x1,\ldots,x50$ (multiples of $|S||A|$) and a right-censored class $>x50$. Our contributions are:
\begin{enumerate}
\item \textbf{Target and benchmark.} A learner-specific, right-censored requirement with a sustained-success rule, and a pipeline that generates controlled MDP families and converts MDPs and behaviors from eight external sources, seven of which enter the final tests, with exact evaluation.
\item \textbf{Method and guarantees.} Offline Effective Horizon (\OEH{}), which simulates the actual learner on MDPs drawn from a posterior fitted to the pilot, adds a planning-risk descriptor and cross-fitted pilot diagnostics, and joins a full and a robust head with a shift gate. We prove that the gate's loss on every forecast is bounded by its components' losses and that the regret of budget decisions is bounded by the forecast's integrated Brier score.
\item \textbf{Pre-registered evaluation.} \OEH{} was developed on opened panels only, frozen, and tested in three protocols on 660 new MDPs (688{,}320 learner jobs), each committed before its data existed, with an architecture-matched control, fixed-sequence gates and MDP-clustered simultaneous bounds. The last protocol registers an endpoint restricted to the budgets the pilot cannot reach and a baseline that trains the learner on pilot prefixes and scores it by fitted Q evaluation (FQE).
\end{enumerate}
\OEH{} beat generic pilot statistics in every final test (registered in Tests 1 and 3, descriptive in Test 2) and the pilot-extrapolation baseline on the extrapolation endpoint. We also report what did not hold: the advantage over the strongest pilot-only control and the benefit of the gate were panel-dependent, and a tabular foundation model on generic statistics matched \OEH{} inside the generator.
